\part{Введение в математический анализ}
\chapter{Теорема Больцано-Вейерштрасса и критерий Коши сходимости числовой последовательности.}

\section{Последовательности и подпоследовательности}

%----------------------------------------------------------------------------------------
% Определение 2.2 --- числовая последовательность
%----------------------------------------------------------------------------------------

\begin{defn} \label{df2:2}
\textit{Числовой последовательностью\rindex{последовательность!числовая}} называется функция с областью определения $\bbN$ и множеством значений, принадлежащим $\bbR$: $y = f(n)$, $n \in \bbN$. Обычно аргумент записывается в виде индекса: $x_n$, $y_n$, $z_n$, $a_n$ и т.д.
\end{defn}

%----------------------------------------------------------------------------------------
% Определение 2.3 --- ограниченность функции и последовательности
%----------------------------------------------------------------------------------------

\begin{defn} \label{df2:3}
Пусть $E(f) \subset \bbR$. Функция $f$ называется \textit{ограниченной\rindex{функция!ограниченная}} (\textit{ограниченной сверху}, \textit{ограниченной снизу}) на множестве $X$, если её множество значений ограничено (ограничено сверху, ограничено снизу). Точная верхняя и нижняя грани $E(f)$ называются \textit{точной верхней} и \textit{нижней гранями} $f$ на $X$ (обозначаются $\sup\limits_{X} f(x)$, $\inf\limits_{X} f(x)$).

Числовая последовательность $x_n$ называется \textit{ограниченной\rindex{последовательность!ограниченная}} (\textit{ограниченной сверху}, \textit{ограниченной снизу}), если множество её значений ограничено (ограничено сверху, ограничено снизу). Точная верхняя и нижняя грани этого множества называются \textit{точной верхней} и \textit{нижней гранями} $x_n$ (обозначаются $\sup x_n$, $\inf x_n$).
\end{defn}

%----------------------------------------------------------------------------------------
% Определение 2.13 --- подпоследовательность
%----------------------------------------------------------------------------------------

\begin{defn} \label{df2:13}
Пусть $x_n$ "--- числовая последовательность, а $n_k$, $k = 1, 2, \ldots$ "--- строго возрастающая последовательность натуральных чисел. Тогда последовательность $y_k = x_{n_k}$ (с индексом $k$) называется \textit{подпоследовательностью\rindex{подпоследовательность}} последовательности $x_n$.
\end{defn}

%----------------------------------------------------------------------------------------
% Определение 2.14 --- частичный предел
%----------------------------------------------------------------------------------------

\begin{defn} \label{df2:14}
Число $a \in \bbR$ называется \textit{частичным пределом\rindex{частичный предел}} (\textit{предельной точкой\rindex{предельная точка}}) последовательности $x_n$, если существует такая строго возрастающая последовательность индексов $n_k$, что $\lim\limits_{k \to \infty} x_{n_k} = a$.
\end{defn}

%----------------------------------------------------------------------------------------
% Теорема 2.3 --- о двух милиционерах
%----------------------------------------------------------------------------------------

\begin{thm}[о двух милиционерах] \label{th2:3}
Если $\lim\limits_{n \to \infty} x_n = \lim\limits_{n \to \infty} z_n = a$ и найдётся номер $n_0$ такой, что при всех $n \ge n_0$ выполнено неравенство $x_n \le y_n \le z_n$, то $\lim\limits_{n \to \infty} y_n = a$.
\end{thm}

%----------------------------------------------------------------------------------------
% Теорема 2.5 --- Кантора о вложенных отрезках
%----------------------------------------------------------------------------------------

\begin{thm}[Кантора о вложенных отрезках] \label{th2:5}
Если $[a_1;\, b_1] \supset [a_2;\, b_2] \supset \ldots \supset [a_n;\, b_n] \supset \ldots$ (бесконечная последовательность вложенных отрезков), то существует точка $\gamma$, общая для всех отрезков (т.е. для всех $n$ выполняется неравенство $a_n \le \gamma \le b_n$). Если при этом последовательность длин отрезков стремится к нулю ($\lim\limits_{n \to \infty} (b_n - a_n) = 0$), то такая точка $\gamma$ единственна, при этом
$$
\gamma = \lim_{n \to \infty} a_n = \sup a_n = \lim_{n \to \infty} b_n = \inf b_n.
$$
\end{thm}

%----------------------------------------------------------------------------------------
% Теорема 2.7 --- критерий частичного предела
%----------------------------------------------------------------------------------------

\begin{thm}[критерий частичного предела] \label{th2:7}
Пусть $\alpha$ "--- один из символов $a$, $+\infty$, $-\infty$. Тогда $\alpha$ является частичным пределом $x_n$ $\iff$ в любой $\varepsilon$-окрестности $\alpha$ ($\varepsilon > 0$) содержится бесконечно много членов $x_n$.
\end{thm}

\begin{notion}
Заметим, что если в любой $U_\varepsilon(\alpha)$ содержится бесконечно много $x_n$, то отсюда ещё не следует, что вне $U_\varepsilon(\alpha)$ не более конечного числа $x_n$ (вне $U_\varepsilon(\alpha)$ тоже может быть бесконечно много $x_n$). Этим и отличается частичный предел от предела последовательности.
\end{notion}

Как мы знаем, ограниченная последовательность может расходиться, но при этом иметь частичные пределы (пример~2.22). Это не случайно, имеет место

%----------------------------------------------------------------------------------------
% Лемма 2.3 --- ограниченность последовательности
%----------------------------------------------------------------------------------------

\begin{lemm} \label{th2:3_b}
Если последовательность $x_n$ ограничена при $n \ge n_0$, $n_0 \in \bbN$ (т.е. $\exists\, m, M \colon \forall\, n \ge n_0 \longrightarrow m \le x_n \le M$), и определена при всех $n \in \bbN$, то она ограничена.
\end{lemm}

%----------------------------------------------------------------------------------------
% Пример 2.22 --- подпоследовательности и частичные пределы
%----------------------------------------------------------------------------------------

\noindent\textbf{Пример 2.22.} Рассмотрим последовательность $x_n = (-1)^n$. Она расходится, но имеет сходящиеся подпоследовательности $x_{2k} = 1$ и $x_{2k-1} = -1$. Таким образом, она имеет частичные пределы $1$ и $-1$.

\section{Теорема Больцано-Вейерштрасса}
%----------------------------------------------------------------------------------------
% Теорема 2.8 --- Больцано–Вейерштрасса
%----------------------------------------------------------------------------------------

\begin{thm}[Больцано"--~Вейерштрасса] \label{th2:8}
Любая ограниченная последовательность имеет сходящуюся подпоследовательность (т.е. имеет конечный частичный предел).
\end{thm}

\begin{proof}
Пусть для всех $n = 1, 2, \ldots$ выполняется неравенство $a \le x_n \le b$. Разобьём отрезок $[a;\, b]$ на 2 равных отрезка $\left[a;\, \dfrac{a+b}{2}\right]$ и $\left[\dfrac{a+b}{2};\, b\right]$; выберем ту половину $\Delta_1$, где содержится бесконечно много членов $x_n$ (и там, и там конечного числа $x_n$ быть не может, так как тогда их всего было бы конечное число). Если и там, и там бесконечно много $x_n$, то $\Delta_1$ "--- любая из половинок.

В отрезке $\Delta_1$ выберем половину $\Delta_2$, где бесконечно много $x_n$ (аналогично), в $\Delta_2$ "--- половину $\Delta_3$, где бесконечно много $x_n$ и т.д. На $k$-м шагу в $\Delta_k$ выберем половину $\Delta_{k+1}$, где бесконечно много $x_n$. Имеем последовательность вложенных отрезков $\Delta_1 \supset \Delta_2 \supset \ldots \supset \Delta_n \supset \ldots$, причём длина $n$-го отрезка равна
$$
\frac{b - a}{2^n} = (b - a) \left(\frac{1}{2}\right)^n \to 0
$$
"--- стремится к нулю.

По теореме Кантора о вложенных отрезках (теорема~\ref{th2:5}) существует единственная точка $c$, принадлежащая всем отрезкам $\Delta_n$. Пусть $\varepsilon > 0$. Так как $\exists\, n_0 \colon \forall\, n \ge n_0 \longrightarrow$ длина $\Delta_n < \varepsilon$, то при $n \ge n_0$ отрезок $\Delta_n$ целиком принадлежит $U_\varepsilon(c)$ (см. рис.~2.10), значит, в $U_\varepsilon(c)$ бесконечно много членов $x_n$. По теореме~\ref{th2:7} $c$ "--- частичный предел $x_n$.
\end{proof}

%----------------------------------------------------------------------------------------
% Теорема 2.9 --- аналог Больцано–Вейерштрасса для неограниченных
%----------------------------------------------------------------------------------------

\begin{thm}[аналог теоремы Больцано"--~Вейерштрасса для неограниченных последовательностей] \label{th2:9}
Если последовательность $x_n$ неограничена сверху, то она имеет частичный предел $+\infty$. Если последовательность $x_n$ неограничена снизу, то она имеет частичный предел $-\infty$.
\end{thm}

\begin{proof}
Докажем первую часть теоремы: вторая доказывается аналогично. Зафиксируем $E > 0$. Так как $x_n$ неограничена сверху, то $\exists\, n_1 \colon x_{n_1} > E$. В качестве нового $E$ в определении неограниченности сверху рассмотрим $x_{n_1}$. Тогда $\exists\, n_2 \colon x_{n_2} > x_{n_1}$. Аналогично, $\exists\, n_3 \colon x_{n_3} > x_{n_2}$ и т.д.

Мы выбрали бесконечно много различных членов последовательности $x_n$ таких, что $x_{n_1} < x_{n_2} < x_{n_3} < \ldots < x_{n_k} < \ldots$ в $U_E(+\infty)$. По теореме~\ref{th2:7} $+\infty$ "--- частичный предел $x_n$.
\end{proof}

\begin{notion}
Итак, любая последовательность имеет частичный предел: ограниченная "--- конечный, неограниченная "--- равный $+\infty$ или $-\infty$.

Отметим, что теорема Больцано"--~Вейерштрасса характерна именно для действительных чисел и выражает свойство их полноты (непрерывности). Её аналог "--- теорема~\ref{th2:9} "--- выполняется и во множестве рациональных чисел.
\end{notion}

%----------------------------------------------------------------------------------------
% Теорема 2.10 --- о единственном частичном пределе
%----------------------------------------------------------------------------------------

\begin{thm}[о единственном частичном пределе] \label{th2:10}
Пусть последовательность $x_n$ ограничена и имеет единственный частичный предел $a$. Тогда последовательность $x_n$ сходится к числу $a$.
\end{thm}


\section{Критерий Коши сходимости числовой последовательности}

%----------------------------------------------------------------------------------------
% Определение 2.17 --- фундаментальная последовательность
%----------------------------------------------------------------------------------------

\begin{defn} \label{df2:17}
Последовательность $x_n$ называется \textit{фундаментальной\rindex{последовательность!фундаментальная}}, если
$$
\forall\, \varepsilon > 0 \;\; \exists\, n_0 \colon \;\; \forall\, n, m \ge n_0 \longrightarrow |x_n - x_m| < \varepsilon
$$
(для любого положительного числа $\varepsilon$ найдётся номер $n_0$ такой, что для любых двух номеров $n \ge n_0$ и $m \ge n_0$ выполняется неравенство $|x_n - x_m| < \varepsilon$).
\end{defn}

%----------------------------------------------------------------------------------------
% Теорема 2.13 --- критерий Коши
%----------------------------------------------------------------------------------------

\begin{thm}[критерий Коши] \label{th2:13}
Последовательность $x_n$ сходится $\iff$ $x_n$ фундаментальна.
\end{thm}

\begin{proof}
\textbf{$\Longrightarrow$} Пусть $\lim\limits_{n \to \infty} x_n = a$. Тогда
$$
\forall\, \varepsilon > 0 \;\; \exists\, n_0 \colon \;\; \left(\forall\, n \ge n_0 \longrightarrow |x_n - a| < \frac{\varepsilon}{2}\right) \wedge \left(\forall\, m \ge n_0 \longrightarrow |x_m - a| < \frac{\varepsilon}{2}\right).
$$
Тогда для любых $n \ge n_0$ и $m \ge n_0$ выполняется неравенство
$$
|x_n - x_m| = |(x_n - a) + (a - x_m)| \le |x_n - a| + |a - x_m| < \frac{\varepsilon}{2} + \frac{\varepsilon}{2} = \varepsilon,
$$
значит, последовательность фундаментальна.

\medskip

\textbf{$\Longleftarrow$} Пусть $x_n$ "--- фундаментальная последовательность. Докажем сначала, что она ограничена. При $\varepsilon = 1$ имеем
$$
\exists\, n_0 \colon \;\; \forall\, n, m \ge n_0 \longrightarrow |x_n - x_m| < 1.
$$
Зафиксируем $m = n_0$. Тогда при $n \ge n_0$ выполнено неравенство
$$
|x_n| = |(x_n - x_{n_0}) + x_{n_0}| \le |x_n - x_{n_0}| + |x_{n_0}| < 1 + |x_{n_0}|.
$$
Таким образом, последовательность $x_n$ ограничена при $n \ge n_0$. По лемме ~\ref{th2:3_b} последовательность ограничена.

По теореме Больцано"--~Вейерштрасса (теорема~\ref{th2:8}) последовательность $x_n$ имеет конечный частичный предел. В силу теоремы~\ref{th2:10} о единственном частичном пределе достаточно доказать, что других частичных пределов последовательность не имеет.

Пусть это не так, и последовательность имеет два различных частичных предела $a$ и $b$ (для определённости $a < b$). Возьмём в определении фундаментальности $\varepsilon = \dfrac{b - a}{3}$ (так, чтобы $U_\varepsilon(a)$ и $U_\varepsilon(b)$ не только не пересекались, но ещё имели между собой зазор ширины $\varepsilon$):
$$
\exists\, n_0 \colon \;\; \forall\, n, m \ge n_0 \longrightarrow |x_n - x_m| < \varepsilon.
$$

Но в $U_\varepsilon(a)$ содержится бесконечно много членов $x_n$ (по теореме~\ref{th2:7}). Значит, $\exists\, n_1 \ge n_0 \colon x_{n_1} \in U_\varepsilon(a)$. Аналогично $\exists\, m_1 \ge n_0 \colon x_{m_1} \in U_\varepsilon(b)$.

Тогда $|x_{n_1} - x_{m_1}| > \varepsilon$. Полученное противоречие показывает единственность частичного предела.
\end{proof}

\begin{notion}
На практике критерий Коши удобно использовать для доказательства расходимости последовательности.
\end{notion}

%----------------------------------------------------------------------------------------
% Пример 2.29 --- применение критерия Коши
%----------------------------------------------------------------------------------------

\noindent\textbf{Пример 2.29.} $x_n = 1 + \frac{1}{2^2} + \dots + \frac{1}{n^2} = \sum\limits_{k=1}^n \frac{1}{k^2}$ (сходимость этой последовательности была установлена при помощи теоремы Вейерштрасса; теперь применим критерий Коши).

\smallskip
$\square$ Имеем
$$
|x_{n+p} - x_n| = \left| \sum_{k=1}^{n+p} \frac{1}{k^2} - \sum_{k=1}^n \frac{1}{k^2} \right| = \sum_{k=n+1}^{n+p} \frac{1}{k^2} < \sum_{k=n+1}^{n+p} \frac{1}{k(k-1)} =
$$
$$
= \sum_{k=n+1}^{n+p} \left( \frac{1}{k-1} - \frac{1}{k} \right) = \frac{1}{n} - \frac{1}{n+p} < \frac{1}{n}.
$$
Это выражение меньше $\varepsilon$ при $n > \frac{1}{\varepsilon}$, т.е. при $n \ge n_0(\varepsilon) = \left[ \frac{1}{\varepsilon} \right] + 1$.

Итак, $\forall\, \varepsilon > 0 \to \exists\, n_0 = \left[ \frac{1}{\varepsilon} \right] + 1 \colon \forall\, n \ge n_0, \forall\, p \in \bbN \to |x_{n+p} - x_n| < \varepsilon$. Последовательность сходится. \hfill$\blacksquare$

\begin{notion}
Отметим, что номер $n_0$ должен зависеть только от $\varepsilon$ и ни в коем случае не должен зависеть от $p$.
\end{notion}

%----------------------------------------------------------------------------------------
% Пример 2.30 --- расходимость гармонической последовательности
%----------------------------------------------------------------------------------------

\noindent\textbf{Пример 2.30.} $x_n = 1 + \frac{1}{2} + \dots + \frac{1}{n} = \sum\limits_{k=1}^n \frac{1}{k}$. Хотя внешне эта последовательность мало отличается от предыдущей, но она расходится.

\smallskip
$\square$ Имеем
$$
|x_{n+p} - x_n| = \sum_{k=n+1}^{n+p} \frac{1}{k} \ge \frac{p}{n+p}
$$
(в сумме $p$ слагаемых, самое маленькое равно $\frac{1}{n+p}$). Возьмём $n = n_0$, $p = n_0$. Тогда $|x_{n+p} - x_n| = \frac{1}{2}$.

Итак, $\exists\, \varepsilon = \frac{1}{2} \colon \forall\, n_0 \to \exists\, n = n_0, \exists\, p = n_0 \colon |x_{n+p} - x_n| = \frac{1}{2}$. Последовательность расходится. \hfill$\blacksquare$

В качестве предостережения приведём неверное «доказательство» того, что эта последовательность сходится.

Имеем $|x_{n+p} - x_n| = \sum\limits_{k=n+1}^{n+p} \frac{1}{k} \le \frac{p}{n+1} < \varepsilon$ при всех $n > \frac{p}{\varepsilon} - 1$. Отсюда нельзя сделать вывод о фундаментальности последовательности $x_n$, так как номер $n_0$ такой, что при $n \ge n_0$ выполняется неравенство $|x_{n+p} - x_n| < \varepsilon$, зависит не только от $\varepsilon$, но и от $p$.

 